% ============================================================
%  Глава 2. Корреляционные деривативы и rainbow-опционы
% ============================================================
\chapter{Корреляционные деривативы и rainbow-опционы}
\label{ch:corr}

\section{Корреляционный своп}
\label{sec:corr-swap}

\subsection{Определение, выплата, экономический смысл}
\label{sec:corr-swap-def}

Корреляционный своп (correlation swap)~--- внебиржевой контракт, по
которому в момент истечения сделки одна сторона получает разницу между
\emph{реализованной} средней парной корреляцией доходностей корзины
активов и заранее зафиксированным \emph{страйком}, умноженную на номинал.
Формально, для корзины из $n$ активов реализованная корреляция за период
$[0,T]$ задаётся как
\begin{equation}
  \rho_{\text{real}}
  \;=\; \frac{2}{n(n-1)}\sum_{1\le i<j\le n}\rho_{ij}^{\text{real}}(0,T),
\end{equation}
где $\rho_{ij}^{\text{real}}(0,T)$~--- парная реализованная корреляция
логдоходностей активов $i$ и $j$ на интервале $[0,T]$.

Контракт характеризуется набором параметров: номинал $N$, срок сделки
$T$, фиксированный страйк $K_{\text{fixed}}$. Выплата покупателю свопа
имеет вид
\begin{equation}
  \label{eq:corrswap-payoff}
  \text{Payoff}\;=\;N\bigl(\rho_{\text{real}} - K_{\text{fixed}}\bigr).
\end{equation}

\placeholderfig{12cm}{6cm}%
{Динамика реализованной корреляции корзины и фиксированного страйка
корреляционного свопа.}{fig:corr-real-vs-strike}

Покупатель корреляционного свопа получает прибыль при росте средней
корреляции корзины и убыток при её падении. Это делает инструмент
естественным для торговли структурой кросс-зависимостей и для
хеджирования портфелей, чувствительных к режимным переходам
коррелированности.

\subsection{Расчёт справедливой стоимости}
\label{sec:corr-models}

В момент заключения сделки её справедливый страйк $K_{\text{fair}}$
определяется условием равенства нулю ожидания выплаты под риск-нейтральной
мерой $\Q$:
\begin{equation}
\label{eq:fair-strike}
  \E^{\Q}[\text{Payoff}] \;=\; N\bigl(\E^{\Q}[\rho_{\text{real}}]
                                   - K_{\text{fixed}}\bigr)\;=\;0
  \quad\Longrightarrow\quad
  K_{\text{fair}} \;=\; \E^{\Q}[\rho_{\text{real}}].
\end{equation}
В непрерывном пределе можно записать
\begin{equation}
\label{eq:rho-real-integral}
  \rho_{\text{real}} \;=\; \frac{1}{T}\int_0^T \rho_t\,\dd t,
  \qquad
  K_{\text{fair}} \;=\; \frac{1}{T}\int_0^T \E^{\Q}[\rho_t]\,\dd t.
\end{equation}
Дальнейший расчёт сводится к выбору модели динамики $\rho_t$ и
вычислению (или приближению) условного среднего $\E^{\Q}[\rho_t]$.

\paragraph{Модель с постоянной корреляцией.}
Простейшее предположение $\rho_t \equiv \rho$ даёт
\begin{equation}
  K_{\text{fair}} \;=\; \rho.
\end{equation}

\paragraph{Трёхпараметрическая модель возврата к среднему.}
Если $\rho_t = \bar\rho + (\rho_0-\bar\rho)e^{-\kappa t}$, то прямое
интегрирование даёт
\begin{equation}
\label{eq:fair-meanrev}
  K_{\text{fair}} \;=\; \bar\rho + (\rho_0-\bar\rho)\frac{1-e^{-\kappa T}}{\kappa T}.
\end{equation}

\paragraph{Модель стохастической корреляции Якоби.}
Стохастическое уравнение
\begin{equation}
\label{eq:jacobi}
  \dd\rho_t \;=\; \kappa(\bar\rho-\rho_t)\,\dd t
                + \sigma\sqrt{1-\rho_t^{\,2}}\,\dd W_t
\end{equation}
обеспечивает $\rho_t \in (-1,1)$ при подходящих параметрах. Из
\eqref{eq:jacobi} следует
\[
  \E^{\Q}[\dd\rho_t] \;=\; \kappa(\bar\rho - \E^{\Q}[\rho_t])\dd t,
\]
поэтому, обозначая $m_t = \E^{\Q}[\rho_t]$, получаем линейное ОДУ
\(\dd m_t/\dd t = \kappa(\bar\rho - m_t)\) с решением
\(
  m_t = \bar\rho + (\rho_0-\bar\rho)e^{-\kappa t}.
\)
Подстановка в~\eqref{eq:rho-real-integral} даёт ту же формулу, что и
в модели возврата к среднему:
\begin{equation}
\label{eq:fair-jacobi}
  K_{\text{fair}} \;=\; \bar\rho
                       + (\rho_0-\bar\rho)\frac{1-e^{-\kappa T}}{\kappa T}.
\end{equation}

\paragraph{Модель стохастической корреляции Фишера--OU.}
Рассмотрим преобразование Фишера $\rho_t = \tanh(x_t)$ с
Орнштейн-Уленбековским процессом
\begin{equation}
\label{eq:fisher-ou}
  \dd x_t \;=\; \kappa(\bar x - x_t)\,\dd t + \eta\,\dd W_t.
\end{equation}
В этом случае
\(
  x_t \;=\; \bar x + (x_0-\bar x)e^{-\kappa t}
            + \eta\!\int_0^t e^{-\kappa(t-s)}\,\dd W_s,
\)
откуда
\[
  m_t \,\equiv\, \E^{\Q}[x_t] \;=\; \bar x + (x_0-\bar x)e^{-\kappa t},
  \qquad
  v_t \,\equiv\, \Var^{\Q}[x_t] \;=\; \frac{\eta^{2}}{2\kappa}\bigl(1-e^{-2\kappa t}\bigr).
\]
Разложение $\tanh$ в окрестности $m_t$ во втором порядке даёт приближение
\begin{equation}
\label{eq:fisher-Erho}
  \E^{\Q}[\rho_t] \;=\; \E^{\Q}[\tanh(x_t)]
  \;\approx\; \tanh(m_t)\bigl(1 - v_t\,\sech^{2}(m_t)\bigr),
\end{equation}
а справедливый страйк~---
\begin{equation}
\label{eq:fair-fisher}
  K_{\text{fair}}
  \;\approx\;
  \frac{1}{T}\int_0^{T}\tanh(m_t)\bigl(1 - v_t\,\sech^{2}(m_t)\bigr)\,\dd t.
\end{equation}

\todoblock{Сюда позже~--- табличка сравнения $K_{\text{fair}}$ по моделям
для типового набора параметров и график зависимости $K_{\text{fair}}(T)$.}
\placeholdertab{Сравнение справедливых страйков $K_{\text{fair}}$ для
рассмотренных моделей при типовом наборе параметров.}{tab:corrmodels}

\subsection{Forward-стоимость и аддитивность по уже реализованным
            значениям}
\label{sec:corr-fv}

Корреляционные свопы с дневной отрисовкой реализованной корреляции
обладают удобным свойством: справедливая стоимость свопа в произвольный
момент $t\in[0,T]$ распадается на детерминированный накопленный $P\&L$
по прошедшему отрезку $[0,t]$ и стоимость <<нового>> свопа на
остаточный срок $[t,T]$.

\begin{proposition}
\label{prop:corr-decomp}
Пусть на $[0,T]$ корреляционный своп с номиналом $N$ имеет страйк
$K_{\text{fixed}}$, а реализованная корреляция на интервале $[0,t]$ уже
известна и равна $\rho^{[0,t]}_{\text{real}}$. Тогда дисконтированная
forward-стоимость свопа в момент $t$ имеет вид
\begin{equation}
\label{eq:corr-decomp}
  V_t \;=\; e^{-r(T-t)}\biggl[
    \tfrac{t}{T}\bigl(\rho^{[0,t]}_{\text{real}}-K_{\text{fixed}}\bigr)
    + \tfrac{T-t}{T}\bigl(K_{\text{fair}}^{[t,T]}-K_{\text{fixed}}\bigr)
  \biggr]\,N,
\end{equation}
где $K_{\text{fair}}^{[t,T]} = \E_t^{\Q}\!\left[\tfrac{1}{T-t}\int_t^T\rho_s\,\dd s\right]$
--- справедливый страйк корреляционного свопа на остаточный срок.
\end{proposition}

Доказательство сводится к разбиению интеграла в~\eqref{eq:rho-real-integral}
на участки $[0,t]$ и $[t,T]$, дисконтированию и переходу к
$\Q$-ожиданию~\cite{shreve2004stoch}. Утверждение~\ref{prop:corr-decomp}
играет ключевую роль при построении динамических хеджей: одну часть
$P\&L$ можно зафиксировать детерминированно, а другую~--- пересчитывать
как стоимость <<свежего>> свопа на остаточный срок.

\placeholderfig{12cm}{6cm}%
{Декомпозиция $V_t$ корреляционного свопа на реализованную и
forward-часть на двух подынтервалах.}{fig:corr-fv-decomp}

\subsection{Хеджирование от пузырей с помощью корреляционных свопов}
\label{sec:corr-hedge-bubble}

Эмпирически кризисные эпизоды на финансовых рынках сопровождаются
\emph{ростом} корреляций между секторами и активами~--- этот эффект
получил название <<correlation breakdown>>.
Кризис 2000~года (dot-com bubble) и кризис 2008~года~--- классические
примеры: падение широкого рынка сопровождается заметным ростом средней
парной корреляции внутри крупных индексов. На этом наблюдении основана
естественная стратегия хеджирования: \emph{покупка} корреляционных свопов
эквивалентна покупке защиты от системных эпизодов одновременного
снижения активов.

Хеджевый портфель на горизонте $T_{\max}$ можно строить как набор
корреляционных свопов разной срочности $\{T_k\}_{k=1}^{K}$ с
номиналами $\{N_k\}$. Номиналы подбираются минимизацией остаточного
$P\&L$ заданной системы базовых активов в стрессовом сценарии:
\begin{equation}
\label{eq:hedge-opt}
  \min_{N_1,\dots,N_K} \;\;
  \E\!\left[\bigl(\Delta V_{\text{port}}^{\text{stress}} + \textstyle\sum_k N_k\,\Delta V_k^{\text{stress}}\bigr)^{\!2}\right].
\end{equation}
Гиперпараметры~--- срочности $T_k$ и сами номиналы $N_k$~--- настраиваются
по историческим стрессовым окнам.

\placeholderfig{12cm}{6cm}%
{Динамика средней корреляции SP500 и P\&L хедж-портфеля из
корреляционных свопов в окнах кризисов 2000 и 2008 годов.}%
{fig:bubble-hedge}

\subsection{Хеджирование rainbow-опционов с помощью корреляционных
           свопов}
\label{sec:corr-hedge-rainbow}

Цена rainbow-опциона зависит, помимо цен и волатильностей базовых
активов, от их парных корреляций. Соответствующая греческая величина~---
<<корреляционная вега>>~--- может быть приближённо хеджирована
покупкой/продажей корреляционных свопов. Формально, разложение цены
rainbow-опциона $C_{\text{rb}}$ в окрестности средней корреляции
$\bar\rho$ имеет вид
\begin{equation}
  C_{\text{rb}}(\rho) \;\approx\; C_{\text{rb}}(\bar\rho)
       + \frac{\partial C_{\text{rb}}}{\partial\rho}\bigg|_{\bar\rho}\,(\rho-\bar\rho),
\end{equation}
и линейная часть точно хеджируется одним корреляционным свопом.
Подробности декомпозиции~--- в~\cref{sec:rainbow-decomp}.

\section{Rainbow-опционы}
\label{sec:rainbow-section}

\subsection{Полу-замкнутая формула и сложность её применения}
\label{sec:rainbow-formula}

Для двух активов в модели Блэка--Шоулза существует полу-аналитическое
выражение для цены rainbow-опционов вида $\max(S^{(1)},S^{(2)})$ и
$\min(S^{(1)},S^{(2)})$ через двухмерную функцию нормального
распределения~\cite{stulz1982options}:
\begin{equation}
\label{eq:rainbow-formula}
  C_{\max}(S^{(1)},S^{(2)};K,T)
  \;=\; S^{(1)}\Phi_2(\cdot) + S^{(2)}\Phi_2(\cdot) - K e^{-rT}\Phi_2(\cdot),
\end{equation}
где аргументы $\Phi_2$ зависят от волатильностей $\sigma_1,\sigma_2$ и
их корреляции $\rho$. Для $n\ge 3$ активов аналитические выражения
становятся неудобны, а необходимость пересчёта при изменении $\rho$
делает прямую формулу плохо приспособленной для практического
динамического хеджирования. Поэтому на практике используется
\emph{декомпозиция} цены.

\subsection{Декомпозиция цены на компоненты}
\label{sec:rainbow-decomp}

Воспользуемся тождеством
\[
  \max(a,b) \;=\; a + (b-a)^{+},
\]
которое для $a=S^{(1)}_T$, $b=S^{(2)}_T$ даёт
\begin{equation}
\label{eq:max-spread}
  \max(S^{(1)}_T,S^{(2)}_T) \;=\; S^{(1)}_T + (S^{(2)}_T-S^{(1)}_T)^{+}.
\end{equation}
Тогда rainbow-call на максимум двух активов разлагается как сумма
форварда на актив~1, спред-опциона между активами~2 и~1 и поправок,
связанных со страйком:
\begin{equation}
\label{eq:rainbow-decomp}
  C_{\text{rb}}^{\max}
  \;=\; \underbrace{C_{\text{vanilla}}(S^{(1)},K,T)}_{\text{закрытая формула}}
  \;+\; \underbrace{C_{\text{spread}}(S^{(2)}-S^{(1)},0,T)}_{\text{полу-замкнутая, чувствительна к }\rho}
  \;+\; \underbrace{R(S^{(1)},S^{(2)},K,T)}_{\text{остаточная нелинейная часть}}.
\end{equation}
\todoblock{Точную декомпозицию вставить из приложенной справки --- сюда
войдут несколько строк аккуратных алгебраических преобразований.}

Идея в том, что первые два слагаемых поддаются аналитической оценке (с
помощью формулы Блэка--Шоулза и формулы Маргрэйба для спред-опциона),
а остаточная часть $R$~--- сравнительно <<гладкая>> функция, которую
удобно аппроксимировать численно или машинно-обучаемой моделью
(\cref{ch:ml}).

\subsection{Динамическое хеджирование}
\label{sec:rainbow-dynhedge}

Динамическое хеджирование rainbow-опциона состоит в построении портфеля
$\Pi_t$ из базовых активов, ванильных опционов и корреляционных свопов,
коэффициенты которых пересчитываются в каждый момент $t$ так, чтобы
производные $\Pi_t$ по всем риск-факторам совпадали с производными
$C_{\text{rb}}$:
\begin{equation}
\label{eq:dyn-hedge}
  \frac{\partial \Pi_t}{\partial S^{(i)}_t} = \frac{\partial C_{\text{rb}}}{\partial S^{(i)}_t},
  \quad
  \frac{\partial \Pi_t}{\partial \sigma_i} = \frac{\partial C_{\text{rb}}}{\partial \sigma_i},
  \quad
  \frac{\partial \Pi_t}{\partial \rho_{ij}} = \frac{\partial C_{\text{rb}}}{\partial \rho_{ij}}.
\end{equation}
Ключевая проблема~--- транзакционные издержки на ребалансировку:
бесконечно частая корректировка экономически невозможна, поэтому на
практике используют дискретную сетку моментов ребалансировки и
учитывают спред/комиссии (см.~\cref{ch:results}).

\placeholderfig{12cm}{6cm}%
{Эволюция P\&L портфеля при динамическом хеджировании
rainbow-опциона с разными частотами ребалансировки.}%
{fig:dyn-hedge-pnl}
